% GENERATED FILE. Edit canonical lessons, then run npm run generate:books.
% canonical-source-hash: fnv1a64:2e220a28adfbf076
% canonical-lessons: BN-C309-ghuri, BN-C309-orano, BN-C309-noboborsho, BN-C309-halkhata, BN-C309-protha

\chapter{Kites and the New Year}
\label{ch:bn-kites-and-the-new-year}
\hlchaptermodality{309}

\begin{chapteropening}
\textbf{By the end of this chapter:} \emph{I can talk about kite flying and the customs of the New Year.}

Complete the last lesson of chapter 309: I can talk about kite flying and the customs of the New Year.
\end{chapteropening}

\section[\texorpdfstring{\bn{ঘুড়ি} (ghuṛi)}{ঘুড়ি (ghuṛi)}]{\bn{ঘুড়ি} (ghuṛi) — a paper kite; a kite (toy)}

\label{lesson:BN-C309-ghuri}



\begin{quote}
Before the new one: say the Bengali for a festive procession, then the Bengali for to embrace.
\end{quote}

\subsection*{You'll want to know: \bn{ঘুড়ি}}
\textbf{\bn{ঘুড়ি}} — \emph{ghuṛi} — \textquotedblleft{}a paper kite; a kite (toy)\textquotedblright{}.

\begin{quote}
\bn{ছুটির} \bn{দিনে} \bn{ছেলেরা} \bn{ছাদে} \bn{ঘুড়ি} \bn{নিয়ে} \bn{খেলছিল।}
\end{quote}

\begin{quote}
\emph{chhuṭir dine chhelerā chhāde ghuṛi niye khelchhilô.}
\end{quote}

\begin{quote}
On the holiday the boys were playing with kites on the roof.
\end{quote}

\begin{cousinweb}
\textbf{\bn{ঘুড়ি}} (\emph{ghuṛi}) is the paper kite; its origin is not settled. A kite the bird is \textbf{\bn{চিল}}.
\end{cousinweb}

\begin{grammarlens}[title={What you've built}]
One more word for telling what happened.
\end{grammarlens}

\subsection*{Your turn}
\begin{itemize}
\raggedright
  \item \emph{Say it:} \emph{ghuṛi}
  \item \emph{Say it:} \emph{ghuṛi}, once more
  \item \emph{Say it:} \emph{chhuṭir dine chhelerā chhāde ghuṛi niye khelchhilô.} — On the holiday the boys were playing with kites on the roof.
  \item \emph{Recall:} say the Bengali for a festive procession, then the Bengali for to embrace, then say \emph{ghuṛi} again
\end{itemize}

\begin{tcolorbox}[breakable,colback=teal!4,colframe=teal!35!black,title={Before you move on}]
What does \bn{ঘুড়ি} mean? (\textquotedblleft{}A paper kite; a kite (toy)\textquotedblright{}.) Say it once more.
\end{tcolorbox}
\section[\texorpdfstring{\bn{ওড়ানো} (oṛāno)}{ওড়ানো (oṛāno)}]{\bn{ওড়ানো} (oṛāno) — to fly a kite; to make fly}

\label{lesson:BN-C309-orano}



\begin{quote}
Before the new one: say the Bengali for to embrace, then the Bengali for a paper kite.
\end{quote}

\subsection*{You'll want to know: \bn{ওড়ানো}}
\textbf{\bn{ওড়ানো}} — \emph{oṛāno} — \textquotedblleft{}to fly a kite; to make fly\textquotedblright{}.

\begin{quote}
\bn{বিকেলে} \bn{বাবা} \bn{আর} \bn{আমি} \bn{ছাদে} \bn{ঘুড়ি} \bn{ওড়ালাম।}
\end{quote}

\begin{quote}
\emph{bikele bābā ār āmi chhāde ghuṛi oṛālām.}
\end{quote}

\begin{quote}
In the late afternoon my father and I flew kites on the roof.
\end{quote}

Past (he/she …): \bn{সেই} \bn{ছেলে} \bn{ঘুড়ি} \bn{ওড়াল}, \emph{sei chhele ghuṛi oṛālô}.

\begin{cousinweb}
\textbf{\bn{ওড়ানো}} (\emph{oṛāno}) is \textbf{\bn{ওড়া}}, to fly — CDIAL 1697, Sanskrit \emph{uḍḍīyatē}, 'flies up' — with the -ānō ending: to make something fly.
\end{cousinweb}

\begin{grammarlens}[title={What you've built}]
One more word for telling what happened.
\end{grammarlens}

\subsection*{Your turn}
\begin{itemize}
\raggedright
  \item \emph{Say it:} \emph{oṛāno}
  \item \emph{Say it:} \emph{oṛāno}, once more
  \item \emph{Say it:} \emph{bikele bābā ār āmi chhāde ghuṛi oṛālām.} — In the late afternoon my father and I flew kites on the roof.
  \item \emph{Recall:} say the Bengali for to embrace, then the Bengali for a paper kite, then say \emph{oṛāno} again
\end{itemize}

\begin{tcolorbox}[breakable,colback=teal!4,colframe=teal!35!black,title={Before you move on}]
What does \bn{ওড়ানো} mean? (\textquotedblleft{}To fly a kite; to make fly\textquotedblright{}.) Say it once more.
\end{tcolorbox}
\section[\texorpdfstring{\bn{নববর্ষ} (nôbôbôrshô)}{নববর্ষ (nôbôbôrshô)}]{\bn{নববর্ষ} (nôbôbôrshô) — the new year}

\label{lesson:BN-C309-noboborsho}



\begin{quote}
Before the new one: say the Bengali for a paper kite, then the Bengali for to fly a kite.
\end{quote}

\subsection*{You'll want to know: \bn{নববর্ষ}}
\textbf{\bn{নববর্ষ}} — \emph{nôbôbôrshô} — \textquotedblleft{}the new year\textquotedblright{}.

\begin{quote}
\bn{নববর্ষের} \bn{দিন} \bn{সবাই} \bn{নতুন} \bn{জামা} \bn{পরে} \bn{বেড়াতে} \bn{যায়।}
\end{quote}

\begin{quote}
\emph{nôbôbôrsher din sôbāi nôtun jāmā pôre beṛāte jāy.}
\end{quote}

\begin{quote}
On the first day of the new year everyone puts on new clothes and goes out visiting.
\end{quote}

\begin{cousinweb}
\textbf{\bn{নববর্ষ}} (\emph{nôbôbôrshô}) joins Sanskrit \emph{nava}, 'new', to \emph{varṣa}, 'rain; a year' (Monier-Williams), the year counted in rainy seasons.
\end{cousinweb}

\begin{grammarlens}[title={What you've built}]
One more word for telling what happened.
\end{grammarlens}

\subsection*{Your turn}
\begin{itemize}
\raggedright
  \item \emph{Say it:} \emph{nôbôbôrshô}
  \item \emph{Say it:} \emph{nôbôbôrshô}, once more
  \item \emph{Say it:} \emph{nôbôbôrsher din sôbāi nôtun jāmā pôre beṛāte jāy.} — On the first day of the new year everyone puts on new clothes and goes out visiting.
  \item \emph{Recall:} say the Bengali for a paper kite, then the Bengali for to fly a kite, then say \emph{nôbôbôrshô} again
\end{itemize}

\begin{tcolorbox}[breakable,colback=teal!4,colframe=teal!35!black,title={Before you move on}]
What does \bn{নববর্ষ} mean? (\textquotedblleft{}The new year\textquotedblright{}.) Say it once more.
\end{tcolorbox}
\section[\texorpdfstring{\bn{হালখাতা} (hālkhātā)}{হালখাতা (hālkhātā)}]{\bn{হালখাতা} (hālkhātā) — a new account book opened at the New Year}

\label{lesson:BN-C309-halkhata}



\begin{quote}
Before the new one: say the Bengali for to fly a kite, then the Bengali for the new year.
\end{quote}

\subsection*{You'll want to know: \bn{হালখাতা}}
\textbf{\bn{হালখাতা}} — \emph{hālkhātā} — \textquotedblleft{}a new account book opened at the New Year\textquotedblright{}.

\begin{quote}
\bn{নববর্ষের} \bn{দিন} \bn{দোকানদাররা} \bn{হালখাতা} \bn{খোলেন।}
\end{quote}

\begin{quote}
\emph{nôbôbôrsher din dokāndārrā hālkhātā kholen.}
\end{quote}

\begin{quote}
On the first day of the new year shopkeepers open a new account book.
\end{quote}

\begin{cousinweb}
\textbf{\bn{হালখাতা}} (\emph{hālkhātā}) joins \bn{হাল}, 'the current state' — usually traced to Arabic \emph{ḥāl} — to \textbf{\bn{খাতা}}, an account book: the fresh ledger a shop opens on the first day of the year, with sweets for its customers.
\end{cousinweb}

\begin{grammarlens}[title={What you've built}]
One more word for telling what happened.
\end{grammarlens}

\subsection*{Your turn}
\begin{itemize}
\raggedright
  \item \emph{Say it:} \emph{hālkhātā}
  \item \emph{Say it:} \emph{hālkhātā}, once more
  \item \emph{Say it:} \emph{nôbôbôrsher din dokāndārrā hālkhātā kholen.} — On the first day of the new year shopkeepers open a new account book.
  \item \emph{Recall:} say the Bengali for to fly a kite, then the Bengali for the new year, then say \emph{hālkhātā} again
\end{itemize}

\begin{tcolorbox}[breakable,colback=teal!4,colframe=teal!35!black,title={Before you move on}]
What does \bn{হালখাতা} mean? (\textquotedblleft{}A new account book opened at the New Year\textquotedblright{}.) Say it once more.
\end{tcolorbox}
\section[\texorpdfstring{\bn{প্রথা} (prôthā)}{প্রথা (prôthā)}]{\bn{প্রথা} (prôthā) — a custom; a tradition}

\label{lesson:BN-C309-protha}



\begin{quote}
Before the new one: say the Bengali for the new year, then the Bengali for a new account book.
\end{quote}

\subsection*{You'll want to know: \bn{প্রথা}}
\textbf{\bn{প্রথা}} — \emph{prôthā} — \textquotedblleft{}a custom; a tradition\textquotedblright{}.

\begin{quote}
\bn{নতুন} \bn{বছরে} \bn{হালখাতা} \bn{খোলা} \bn{একটা} \bn{পুরনো} \bn{প্রথা।}
\end{quote}

\begin{quote}
\emph{nôtun bôchhôre hālkhātā kholā ekṭā purôno prôthā.}
\end{quote}

\begin{quote}
Opening a new ledger for the new year is an old custom.
\end{quote}

\begin{cousinweb}
\textbf{\bn{প্রথা}} (\emph{prôthā}) is Sanskrit \emph{prathā}, 'a spreading out; fame' (Monier-Williams); the sense 'a custom', a practice that has spread, is the later one.
\end{cousinweb}

\begin{grammarlens}[title={What you've built}]
One more word for telling what happened.
\end{grammarlens}

\subsection*{Your turn}
\begin{itemize}
\raggedright
  \item \emph{Say it:} \emph{prôthā}
  \item \emph{Say it:} \emph{prôthā}, once more
  \item \emph{Say it:} \emph{nôtun bôchhôre hālkhātā kholā ekṭā purôno prôthā.} — Opening a new ledger for the new year is an old custom.
  \item \emph{Recall:} say the Bengali for the new year, then the Bengali for a new account book, then say \emph{prôthā} again
\end{itemize}

\begin{tcolorbox}[breakable,colback=teal!4,colframe=teal!35!black,title={Before you move on}]
What does \bn{প্রথা} mean? (\textquotedblleft{}A custom; a tradition\textquotedblright{}.) Say it once more.
\end{tcolorbox}
